\documentclass[final]{beamer}

% Poster layout (beamerposter) — portrait, two columns
\usepackage[size=a0,scale=1.08,orientation=portrait]{beamerposter}
\usetheme{default}
\usecolortheme{default}

\usepackage[T1]{fontenc}
\usepackage[utf8]{inputenc}
\usepackage[english]{babel}
\usepackage{lmodern}
\usepackage{microtype}
\usepackage{amsmath,amssymb,mathtools}
\usepackage{booktabs}
\usepackage{hyperref}
\usepackage{graphicx}
\usepackage{xcolor}

% --- Dark blue color scheme ---
\definecolor{darkblue}{RGB}{0,51,102}
\setbeamercolor{block title}{fg=white,bg=darkblue}
\setbeamercolor{block body}{fg=black,bg=white}
\setbeamerfont{block title}{series=\bfseries,size=\Large}
\setbeamerfont{block body}{size=\large}
\setbeamertemplate{block begin}{
  \vskip0.75ex
  \begin{beamercolorbox}[wd=\textwidth,sep=1.2ex,rounded=true,leftskip=0.5ex,rightskip=0.5ex]{block title}
    \usebeamerfont{block title}\insertblocktitle
  \end{beamercolorbox}
  \begin{beamercolorbox}[wd=\textwidth,sep=1.2ex,rounded=true,leftskip=0.5ex,rightskip=0.5ex]{block body}
    \usebeamerfont{block body}
}
\setbeamertemplate{block end}{
  \end{beamercolorbox}
  \vskip0.75ex
}

% Frame style for ideas
\newcommand{\framedidea}[1]{%
  \vspace{0.3em}
  \setlength{\fboxrule}{2pt}
  \setlength{\fboxsep}{0.8ex}
  \fcolorbox{darkblue}{white}{%
    \parbox{0.92\textwidth}{#1}%
  }%
  \vspace{0.3em}
}

% Header
\title{Minimum Sizes of Identifying Codes in Graphs Differing by One Vertex (OH2)}
\author{Charon, Honkala, Hudry, Lobstein (2013), \emph{Cryptography and Communications}}
\institute{Poster author: Vielzeuf Charles}
\date{}

\begin{document}
\begin{frame}[t]

% --- Big title banner with dark blue ---
\begin{beamercolorbox}[wd=\textwidth,sep=1.6ex,center,rounded=true]{block title}
  \huge\bfseries \inserttitle\\[0.5em]
  \Large \insertauthor \qquad\textbar\qquad \large \insertinstitute
\end{beamercolorbox}

\vspace{1.2em}

\begin{columns}[t,totalwidth=\textwidth]

% ===================== Column 1 (up to and including Lemmas) =====================
\begin{column}{0.48\textwidth}

\begin{block}{Context and Motivation}
\textbf{Identifying codes} model the placement of sensors on certain vertices of a graph to \textbf{identify} (or locate) an unknown vertex from the sensors that ``see'' it.

\framedidea{%
\textbf{Example application: Museum rooms}\\
Consider a museum with rooms connected by corridors. We want to place detectors (sensors) in some rooms so that if an alarm is
triggered, we can identify exactly which room it came from based on which detectors detect it. The detectors have a detection radius $r$:
they can detect alarms in rooms within distance $r$. The goal is to minimize the number of detectors needed while ensuring every room can be uniquely identified.
We note by \textbf{$\gamma_r(G)$ the minimum number of detectors needed to identify all rooms in a graph $G$}.
}

Article OH2 studies the \textbf{stability} of the parameter $\gamma_r(G)$ when modifying the graph by \textbf{adding or deleting a vertex}.
With $G^*$ obtained from $G$ by adding/deleting a vertex, provided that $G$ and $G^*$ remain \textbf{identifiable}, we compare:
$$
\gamma_r(G^*)-\gamma_r(G)
\qquad\text{and}\qquad
\frac{\gamma_r(G^*)}{\gamma_r(G)}.
$$
\end{block}

\begin{block}{Main Definitions}
Let $G=(V,E)$ be a simple undirected graph. For $v\in V$ and $r\ge 1$, the \textbf{ball} of radius $r$ is
$ B_{G,r}(v)=\{x\in V: d_G(v,x)\le r\}. $

\framedidea{%
Two vertices $x\ne y$ are \textbf{$(G,S,r)$-twins} if
$$
B_{G,r}(x)\cap S = B_{G,r}(y)\cap S.
$$
The graph is \textbf{$r$-twin-free} if no such twins exist (when $S=V$).
}

A \textbf{code} is a subset $\mathcal{C}\subseteq V$; the elements of C are called \textbf{codewords}.

For a vertex $v$, the associated \textbf{identifying set} is:
$$
I_{G,\mathcal{C},r}(v)=B_{G,r}(v)\cap \mathcal{C}.
$$

\framedidea{%
The code $\mathcal{C}$ is \textbf{$r$-identifying} if the sets $I_{G,\mathcal{C},r}(v)$, $v\in V$, are:
\begin{itemize}
  \item \textbf{nonempty} (every vertex is covered),
  \item \textbf{all distinct} (every pair of vertices is separated).
\end{itemize}
}
A graph admits an $r$-identifying code only if it is $r$-twin-free: two $r$-twins would have the same identifying set $I_{G,\mathcal{C},r}(\cdot)$ for any code $\mathcal{C}$, so they could not be separated.

When $G$ is $r$-twin-free, we define
$
\gamma_r(G)=\min\{|\mathcal{C}|:\ \mathcal{C}\ \text{is an $r$-identifying code of}\ G\}.
$

\vspace{1.5em}

\begin{center}
\includegraphics[width=\textwidth]{cases.png}
\end{center}
\end{block}

\begin{block}{Main Questions}
\framedidea{%
Given an $r$-twin-free graph $G$, and $G^*$ obtained by \textbf{adding} or \textbf{deleting} a vertex,\\
if $G^*$ is still $r$-twin-free:
\begin{itemize}
  \item What differences are possible for $\gamma_r(G^*)-\gamma_r(G)$?
  \item What ratios are possible for $\gamma_r(G^*)/\gamma_r(G)$?
\end{itemize}
}

OH2 provides \textbf{extremal constructions} showing that these quantities can vary strongly, including for \textbf{connected} graphs.
\end{block}

\begin{block}{Lemmas}
\noindent\textbf{Lemma 1} [(G, S, r)-twin transitivity] In a graph $G = (V, E)$, if $x, y, z$ are three distinct vertices, if $S$ is a subset of $V$, if $x$ and $y$ are $(G, S, r)$-twins and if $y$ and $z$ are $(G, S, r)$-twins, then $x$ and $z$ are $(G, S, r)$-twins.


\hspace*{0.8em}\textbf{Lemma 2} If $C$ is an $r$-identifying code in a graph $G = (V, E)$, then so is any set $S$ such that $C \subseteq S \subseteq V$.
\end{block}

\begin{block}{Case $r=1$}
\framedidea{%
Let $k \geq 1$ be an arbitrary integer. There exist two (connected) $1$-twin-free graphs $G$ and $G_x$, where $G_x$ denotes the graph obtained from $G$ by deleting the vertex $x$ and $G$ has $2k + \lceil\log_2(k + 1)\rceil + 2$ vertices, such that
$$
\gamma_1(G) \leq \lceil\log_2(k + 1)\rceil + 2 \quad \text{and} \quad \gamma_1(G_x) \geq k.
$$
}
\end{block}

\end{column}

% ===================== Column 2 (Proof Case r=1, Case r>=2, Conclusion) =====================
\begin{column}{0.48\textwidth}

\textbf{Proof:} Construction with two similar sets of $x$- and $y$-vertices, each identifiable via the codewords $a$, and a vertex $x$ adjacent to only one of these sets.

\vspace{0.8em}

\begin{center}
\includegraphics[width=0.72\linewidth]{example_graph_proof.png}\\
\small Example $k=7$: the codewords $a_i$ + $x$ suffice to identify all $x_i$ and $y_i$ in $G$. $x$ is the vertex removed to get $G_x$;
\end{center}

\textbf{Open question:} Can we do better?

\vspace{1em}

\framedidea{%
Let $G = (V, E)$ be any $1$-twin-free graph with at least three vertices. For any vertex $x \in V$ such that $G_x$ is $1$-twin-free, we have:
$$
\gamma_1(G_x) \geq \gamma_1(G) - 1.
$$
}

\textbf{Proof:} From a minimum $1$-identifying code in $G_x$, one builds a $1$-identifying code in $G$ of size at most $\gamma_1(G_x)+1$ by adding $x$ if it is not $1$-covered, or (when $x$ is $1$-covered but $C - \{x\}$ fails to $1$-separate $x$ from some vertex $v$) by adding a vertex $z$ that $1$-covers exactly one of $x$ and $v$ (such $z$ exists because $G$ is $1$-twin-free); hence $\gamma_1(G)\le\gamma_1(G_x)+1$.

\vspace{1em}

\textbf{Corollary:} If $\gamma_1(G_x) \leq a$ and $\gamma_1(G) \geq a + 1$, then $\gamma_1(G_x) = a$ and $\gamma_1(G) = a + 1$.

\vspace{1em}

\begin{block}{Case $r \geq 2$}
\framedidea{%
For $r \geq 2$, both $\gamma_r(G_x) - \gamma_r(G)$ and $\gamma_r(G_x)/\gamma_r(G)$ can be \textbf{arbitrarily large or small}.
}

\textbf{General view.} The paper uses cycles (e.g.\ $\gamma_r(C_n) = n/2$ for even $n \geq 2r+4$), then connected graphs with \textit{$r$ even} and \textit{$r$ odd} where the difference grows with $n$; path-based constructions yield ratios that can be made arbitrarily large; other families give exact $\gamma_r$ for graphs of order $pr+1$ and $pr+2$. Both the difference and the ratio can be made arbitrarily large or arbitrarily small (see best know ratios from conclusion).
\end{block}

\begin{block}{General Conclusion}
Table~1 summarizes the results of the paper: possible ranges for the difference $\gamma_r(G_x) - \gamma_r(G)$ and the ratio $\gamma_r(G_x)/\gamma_r(G)$ when $G_x$ is obtained from $G$ by adding or deleting one vertex. The asymptotic behaviour is given for $n$ large with respect to $r$; $\approx$ denotes approximate worth known growth.

\vspace{0.5em}

\begin{center}
\large
\begin{tabular}{cclcc}
\toprule
$r$ & $r$ & comment & $\gamma_r(G_x) - \gamma_r(G)$ & $\gamma_r(G_x) / \gamma_r(G)$ \\
\midrule
$r = 1$ & & always & $\geq -1$ & \\
$r = 1$ & & (connected) graphs & $\gtrsim 0.5 n - 1.5 \log_2 n$ & $\gtrsim n / (2 \log_2 n)$ \\
$\geq 2$ & even & connected graphs & $\gtrsim n/4$ & \\
$\geq 2$ & odd & connected graphs & $\gtrsim n(3r-1) / (12r)$ & \\
$\geq 2$ & any & not connected graphs & $\gtrsim n(2r-2) / (2r+1)$ & \\
$\geq 2$ & any & (connected) graphs & & $\gtrsim n / (2r^2 \log_2 n)$ \\
$\geq 2$ & any & (connected) graphs & $\lesssim -n(r-1) / r$ & $\lesssim r(r+1) \log_2 n / n$ \\
\bottomrule
\end{tabular}
\end{center}

\vspace{0.3em}
\small
\textbf{Table 1:} The difference $\gamma_r(G_x) - \gamma_r(G)$ and ratio $\gamma_r(G_x) / \gamma_r(G)$, as functions of $n$ and $r$.
\end{block}

\begin{block}{Sources}
\textbf{More results}
\begin{itemize}
  \item A. Lobstein, O. Hudry, I. Charon, «~Locating-domination and identification~», Ch.~6 in \emph{Topics in Domination in Graphs}, eds.~T.~Haynes, S.~Hedetniemi, M.~Henning, Springer, 2020, pp.~251--299.
  \item O. Hudry, V. Junnila, A. Lobstein, «~On Iiro Honkala's contributions to identifying codes~», \emph{Fundamenta Informaticae} \textbf{191} (3--4), 2024, pp.~165--196.
\end{itemize}

\vspace{0.6em}
\textbf{Applications of identifying codes}
\begin{itemize}
  \item K. Basu, A. Sen, «~Identifying individuals associated with organized criminal networks: A social network analysis~», \emph{Social Networks} \textbf{64}, 2021, pp.~42--54.
  \item K. Basu, M. Padhee, S. Roy, A. Pal, A. Sen, M. Rhodes, B. Keel, «~Health Monitoring of Critical Power System Equipments using Identifying Codes~», preprint, 2018.
  \item S. Sengupta, K. Basu, A. Sen, S. Kambhampati, «~Moving Target Defense for Robust Monitoring of Electric Grid Transformers in Adversarial Environments~», \emph{GameSec} 2020. Uses \textbf{discriminating codes} (variant of identifying codes).
\end{itemize}
\end{block}

\end{column}

\end{columns}

\end{frame}
\end{document}
